\section{怎样识别和构造凸函数}
\label{sec:09-Convex-Optimization/03-Convex-Functions/04}

你正在搭一个损失函数。log-sum-exp 包着线性预测，外面还有正则项，最后对样本求和。你隐约感觉这应该是凸的——但如果必须回到弦定义、对随机选的两点手工验证不等式，这一关就过不去。书里列一排凸函数只是起点；真正的功夫在于知道哪些操作可以反复叠加而保持凸性。

本节整理一套保凸的合法操作。每条规则都经过证明，你的工作是在建模时快速判断"这样写安不安全"，而非每次都从第一原理出发重走一遍。

\subsection{非负加权和与逐点上确界}

最安全的两步。

若每个 \(f_i\) 都凸，且 \(\alpha_i\ge0\)，那么

\[
f(x)=\sum_i\alpha_i f_i(x)
\]

仍凸。弦不等式逐项相加，权重非负保证不等号方向不变。

经验风险

\[
\frac1n\sum_{i=1}^n\ell_i(w)+\lambda R(w)
\]

因此可以通过凸损失加凸正则构造。这是机器学习中最日常的操作，也是 DCP 规则库里的第一条。

凸函数族的逐点上确界

\[
f(x)=\sup_{\alpha} f_\alpha(x)
\]

也是凸函数。理由是上境图的交仍凸——\hyperref[sec:09-Convex-Optimization/03-Convex-Functions/01]{``弦在图像上方意味着什么''}一节对逐点最大值用的正是这个论证。这里不只是一个数学事实；它解释了大量凸模型的结构来源：

\begin{itemize}
\tightlist
\item
  分段线性函数 = 若干仿射函数的最大值：每个仿射函数是一条折线的"面板"，取最大值就是把它们拼在一起；
\item
  凸集 \(C\) 的支持函数 \(\sigma_C(y)=\sup_{x\in C}y^{\trans}x\)：把内积扫过整个集合；
\item
  对称矩阵的最大特征值 \(\lambda_{\max}(X)=\sup_{\norm{v}=1}v^{\trans}Xv\)：Rayleigh 商的上确界。
\end{itemize}

看到 sup，往凸的方向想一步；这条习惯能省大量 Hessian 验算。

逐点下确界则危险得多。两只碗取较低的那只，中间可能出现一个凹陷的鞍部——不再是凸的。但有一个重要例外：如果对一部分变量取最小化，且联合函数凸，那么

\[
g(x)=\inf_y f(x,y),
\]

在适当定义域条件下 \(g\) 仍凸。逻辑是联合上境图投影到 \(x\) 空间后仍凸。消去辅助变量、构造 value function 时，这条规则频繁登场。

\subsection{仿射复合总是安全的}

若 \(f\) 凸，则

\[
g(x)=f(Ax+b)
\]

凸。仿射映射忠实保存凸组合：

\[
A(\theta x+(1-\theta)y)+b = \theta(Ax+b)+(1-\theta)(Ay+b).
\]

弦条件原封不动传递过去，不需要附加条件。最小二乘 \(\norm{ Ax-b}_2^2\)、范数残差 \(\norm{ Ax-b}\) 和 logistic 损失中的线性预测项，都依赖这条最简单的规则。

一般复合则需要单调性配合。若外层 \(h\) 凸且逐坐标非减，内层 \(g_i\) 凸，那么

\[
h(g_1(x),\ldots,g_k(x))
\]

凸。单调方向是这里的关键：外层非减 \(\rightarrow\) 内层需要凸；外层非增 \(\rightarrow\) 对应的内层需要凹。方向错一个，整个复合的凸性就不能保证。

例如 \(e^{g(x)}\) 在 \(g\) 凸时仍凸——指数既凸又在全实轴上递增，两个条件都满足。但 \(g(x)^2\) 不能仅因平方函数凸就断言凸：平方在负半轴上递减，单调方向不匹配。如果你能额外确认 \(g(x)\ge0\)，就把定义域限制在了平方的递增区域，这条复合才安全。DCP 系统在检查复合时正是这样同时追踪曲率和单调性标记的。

\subsection{Log-sum-exp：从稳定计算到凸模型}

\[
f(x)=\log\sum_{i=1}^n e^{x_i}
\]

这个函数在数值分析中出现是因为它安全地计算 log-sum-exp 而不会溢出；在凸分析中，它则是最大值的光滑近似——当某个 \(x_i\) 远大于其他项时，log-sum-exp 趋近于 \(\max_i x_i\)。

其梯度是 softmax 概率向量 \(p\)，即 \(p_i=e^{x_i}/\sum_j e^{x_j}\)，Hessian 为

\[
\nabla^2 f(x) = \diag(p)-pp^{\trans}.
\]

对任意方向 \(v\)，

\[
v^{\trans}\nabla^2f(x)v = \sum_i p_i v_i^2-\Bigl(\sum_i p_i v_i\Bigr)^2
=\Var_p(v_i)\ge0.
\]

右端就是"以概率 \(p_i\) 取值 \(v_i\)"的随机变量的方差——这里的等号是恒等式，不是类比。方差非负，所以 Hessian 半正定，LSE 凸。这个证明还顺便揭示：沿全 \(\mathbf{1}\) 方向 Hessian 为零。因为给所有 \(x_i\) 加同一常数只会让函数值整体平移该常数，softmax 不变——你希望光滑近似具备的这个平移不变性，正好写在零特征值的方向里。

二分类 logistic 损失

\[
\ell(t)=\log(1+e^{-t})
\]

是 LSE 与仿射映射的组合：\(\log(e^0+e^{-t})\)。因此凸。代入 \(t=y_i w^{\trans}x_i\) 后对样本求和，整个经验风险仍凸。从一个原子加两条规则推导整个损失函数的凸性——这套流程比每次重新算 Hessian 干净得多。

\subsection{透视函数把比例关系变成凸结构}

若 \(f\) 凸，其透视函数定义为

\[
g(x,t)=t f(x/t),
\qquad t>0.
\]

透视保持凸性。以 \(f(x)=x^2\) 为例，

\[
g(x,t)=\frac{x^2}{t},
\qquad t>0,
\]

从分式外观看可能觉得不凸——毕竟分母出现了变量。但它确实是凸函数。二次除线性、若干分式模型和锥约束都来自透视函数。

证明可以从上境图的线性变换入手。实用上，重要的是认出"同质化"结构：如果表达式能写成 \(t\cdot f(x/t)\) 的形式且 \(f\) 凸，那它就没问题。不要因为分母出现了变量就立刻判定为非凸。

\subsection{保凸规则是一种建模纪律}

Disciplined Convex Programming（DCP）要求表达式由已知曲率与单调性的原子按合法规则组合。求解器在运行前据此验证凸性并转换为标准锥形式；它检查合法性，但不替你判断建模是否合理。

这套纪律有意做得保守：有些数学上确定凸的函数，不能被现有规则直接识别。遇到这种情况，你可以寻找等价的 DCP 表达、补充新原子，或给出独立证明后关闭该表达式的检查——而不是关掉全局验证，让求解器在黑箱中碰运气。

同样，以下外观判断都不可靠——每条都对应一类凭直觉可能踩的坑：

\begin{itemize}
\tightlist
\item
  两个正凸函数的乘积不一定凸：\(x\cdot e^x\)（\(x>0\)）是凸的吗？去验证。
\item
  两个凸函数的差不一定凸：\(x^2-|x|\) 在原点附近的行为就出问题。
\item
  凸函数的等值约束一般不凸：\(x^2+y^2=1\) 是单位圆周，连线穿过内部。
\item
  可逆的非线性变量变换不一定保凸：换元后的函数可能在新的坐标里弯向错误的方向。
\item
  Hessian 在若干采样点半正定不能证明全域凸：局部信息和全局性质之间没有捷径。
\end{itemize}

所有这些坑的背后是同一个教训：凸性是全局性质，局部外观和局部检验都不足以担保。一条规则的安全边界写在它的证明条件里，超出边界就需要新的论证。
